%% Drafted from the mapped corpus by scripts/draft_topics.py.
% Model: gpt-5.6-luna; source characters: 14340.
\section{Бази от данни}

\subsection{Предназначение на базите от данни}

База от данни е организирана колекция от данни, предназначена за съхраняване, обработване и извличане на информация. Данните се описват чрез определена структура, а операциите върху тях се изпълняват посредством специализиран език за работа с бази от данни.

Релационната база от данни представлява колекция от релации. В практическа реализация релациите обикновено се представят като таблици, но математическото им разглеждане се основава на множества от кортежи. Релационният модел дава възможност данните да бъдат описани чрез формална схема и да бъдат обработвани чрез операции от релационната алгебра или чрез заявки на SQL.

Проектирането на релационна база от данни включва следните основни стъпки:

\begin{enumerate}
    \item определяне на данните, които ще се съхраняват;
    \item определяне на връзките между обектите;
    \item определяне на ключовете и свързващите колони;
    \item определяне на ограниченията върху обектите и връзките между тях;
    \item отстраняване на евентуални недостатъци и излишества;
    \item реализиране на базата от данни.
\end{enumerate}

Тези стъпки водят от описание на предметната област към конкретна структура от релации и атрибути.

\section{Релационен модел на данните}

\subsection{Математическа основа}

Релационният модел се основава на математическото понятие за $n$-членна релация. Една $n$-членна релация е множество от елементи, всеки от които се състои от $n$ компоненти. Тези елементи се наричат $n$-торки или кортежи.

Чрез една релация се моделира даден клас от обекти, а всеки кортеж от релацията представя конкретен обект от този клас. Например релация за филми може да описва класа от всички филми, като всеки неин кортеж представя един конкретен филм.

\begin{defn}
Домейн е именувано множество, разглеждано като множество от допустими стойности на дадена величина.
\end{defn}

Домейнът определя какви стойности могат да се използват за даден атрибут. Например атрибутът \texttt{title} може да има домейн от низове, а атрибутът \texttt{year} може да има домейн от цели числа:

\[
\texttt{title}:\texttt{string}, \qquad
\texttt{year}:\texttt{integer}.
\]

\begin{defn}
В контекста на базите от данни релация е таблица, чиито редове са кортежи, а стойностите в една колона принадлежат на домейна, асоцииран със съответния атрибут.
\end{defn}

Релацията е множество от уникални кортежи. Следователно в математическия релационен модел един и същ кортеж не се повтаря.

\begin{defn}
Атрибут е име на колона в релацията, с което е свързан определен домейн.
\end{defn}

Например релацията за филми може да има атрибути \texttt{title}, \texttt{year}, \texttt{length} и \texttt{filmType}. Съответните домейни могат да бъдат низове, цели числа и други прости типове.

\begin{defn}
Кортеж е ред от релацията, съдържащ конкретна стойност за всеки атрибут.
\end{defn}

Ако атрибутите са подредени като
\[
(\texttt{title},\texttt{year},\texttt{length},\texttt{filmType}),
\]
примерен кортеж е
\[
(\texttt{Star\ wars},1977,124,\texttt{Color}).
\]

Компонентите на всеки кортеж принадлежат на домейните на съответните атрибути. Кортежите следват предварително определената подредба на атрибутите от схемата на релацията.

Релационният модел изисква стойностите на компонентите да бъдат атомарни. Това означава, че във всяка позиция на кортежа се съдържа една проста стойност, например цяло число или низ. Сложни стойности като списъци и масиви не се разглеждат като един компонент на кортежа в основния релационен модел.

\subsection{Схема и екземпляр на релация}

\begin{defn}
Схема на релация е името на релацията заедно с множеството от нейните атрибути.
\end{defn}

Схемата се записва чрез името на релацията, последвано от имената на атрибутите в скоби:

\[
\textit{Movies}(\textit{title},\textit{year},\textit{length},\textit{filmType}).
\]

В математическото описание атрибутите образуват множество. При записването на схемата обаче се фиксира стандартна подредба, за да бъде еднозначно определена позицията на всеки компонент в кортежа.

\begin{defn}
Екземпляр на релация е множеството от кортежи, които в даден момент се съдържат в релацията.
\end{defn}

Схемата описва структурата на релацията, докато екземплярът описва нейното текущо съдържание. Схемата обикновено се променя по-рядко, а екземплярът може да се изменя при добавяне, промяна или изтриване на данни.

\begin{defn}
Кардиналност на екземпляр на релация е броят на кортежите в този екземпляр.
\end{defn}

Ако релацията \textit{Movies} съдържа $120$ кортежа, нейната кардиналност за съответния момент е $120$.

\begin{defn}
Схема на релационна база от данни е съвкупността от схемите на всички релации в базата от данни.
\end{defn}

Например една база от данни може да има следната схема:

\[
\begin{aligned}
\textit{Movie}(&\textit{Title},\textit{Year},\textit{Length}),\\
\textit{StarsIn}(&\textit{MovieTitle},\textit{MovieYear},\textit{StarName}).
\end{aligned}
\]

Подчертаването се използва, за да се означат атрибути, които в дадения пример участват в идентифицирането на кортежите. Така в \textit{Movie} заглавието, годината и дължината са показани като определящи атрибути, а в \textit{StarsIn} са показани \textit{MovieTitle} и \textit{StarName}. Връзката между двете релации се изразява чрез съответните атрибути за филм и година.

\section{Операции върху релационна база от данни}

Операциите върху релационната база от данни могат да бъдат разгледани според това дали описват структурата, обработват данните или управляват достъпа и запазването на промените.

\subsection{DDL операции}

Data Definition Language, или DDL, е подезик за описание и промяна на схемата на базата от данни. Основните DDL команди са \texttt{CREATE}, \texttt{ALTER} и \texttt{DROP}.

Пример за създаване на таблица е:

\begin{verbatim}
CREATE TABLE movies (
  id INTEGER PRIMARY KEY,
  title CHAR(50) NOT NULL,
  length INTEGER NULL
);
\end{verbatim}

Тук се създава таблица \texttt{movies} с три колони. Колоната \texttt{id} е от тип \texttt{INTEGER} и е обявена като първичен ключ. Колоната \texttt{title} съдържа низ с дължина до 50 символа и е обявена като задължителна чрез \texttt{NOT NULL}. За \texttt{length} е разрешена празна стойност чрез \texttt{NULL}.

Промяна на схемата може да се извърши чрез \texttt{ALTER TABLE}:

\begin{verbatim}
ALTER TABLE movies ADD release_date DATE;
ALTER TABLE movies DROP COLUMN length;
\end{verbatim}

Първата команда добавя атрибута \texttt{release\_date}, а втората премахва атрибута \texttt{length}. Изтриването на цялата таблица се извършва чрез:

\begin{verbatim}
DROP TABLE movies;
\end{verbatim}

При изпълнение на DDL операция промените върху схемата се прилагат незабавно.

\subsection{DML операции}

Data Manipulation Language, или DML, е подезик за работа със съдържанието на релациите. Основните DML команди са \texttt{SELECT}, \texttt{INSERT}, \texttt{UPDATE} и \texttt{DELETE}.

Командата \texttt{SELECT} извлича данни:

\begin{verbatim}
SELECT title FROM movies
WHERE length IS NOT NULL AND length < 120
ORDER BY release_date DESC
\end{verbatim}

Заявката извлича заглавията на филмите, чиято дължина не е празна и е по-малка от $120$. Резултатът се подрежда по \texttt{release\_date} в низходящ ред.

Добавяне на кортеж се извършва с \texttt{INSERT}:

\begin{verbatim}
INSERT INTO movies (title, length, release_date)
VALUES ('Star wars', 124, '09-04-1984')
\end{verbatim}

Промяна на вече съществуващи данни се извършва с \texttt{UPDATE}:

\begin{verbatim}
UPDATE movies
SET title = 'Star Wars'
WHERE title = 'Star wars'
\end{verbatim}

Изтриване на кортежи се извършва чрез \texttt{DELETE}:

\begin{verbatim}
DELETE FROM movies
WHERE length < 120
\end{verbatim}

След изпълнение на DML операция промените трябва да бъдат потвърдени с \texttt{COMMIT}, за да станат постоянни.

\subsection{Заявки към релационната база от данни}

Заявка е описание на желан резултат, получен чрез обработване на една или повече релации. В SQL заявките най-често се използва командата \texttt{SELECT}, която може да включва избиране на колони, условие за подбор на редове и подреждане на резултата.

В релационната алгебра заявката се представя чрез алгебричен израз. Всеки оператор приема релация или релации и връща нова релация. По този начин сложните заявки могат да се изграждат чрез комбиниране на по-прости операции.

\section{Релационна алгебра}

\begin{defn}
Релационната алгебра е междинен език за изчисляване на заявки, който задава правила за обработване на релации чрез алгебрични изрази.
\end{defn}

Операциите в релационната алгебра се изпълняват върху релации. В ядрото на релационната алгебра релациите се разглеждат като множества от кортежи. Разширената релационна алгебра разглежда релациите като мултимножества от кортежи, при които могат да се срещат повторения.

Основните класове операции са:

\begin{itemize}
    \item операции върху множества;
    \item операции за отсяване на части от релациите;
    \item операции за комбиниране на кортежи от две релации;
    \item операция за преименуване.
\end{itemize}

Основните операции са селекция, проекция, обединение, разлика, декартово произведение и преименуване. Сечение, тета-съединение и естествено съединение могат да се изразят чрез основните операции и поради това се разглеждат като допълнителни.

\subsection{Обединение}

Обединението на две релации $R$ и $S$ се означава с

\[
R\cup S.
\]

Операцията е приложима, когато релациите имат съвместими схеми: еднакъв брой атрибути, които си съответстват. Резултатът е релация със същата схема, съдържаща всички кортежи, които принадлежат на $R$ или на $S$.

Обединението е бинарна операция, комутативна и асоциативна:

\[
R\cup S=S\cup R,
\]

\[
(R\cup S)\cup T=R\cup(S\cup T).
\]

Повтарящите се кортежи не се включват многократно, тъй като релациите в ядрото на релационната алгебра са множества.

\subsection{Разлика}

Разликата на $R$ и $S$ се означава с

\[
R-S.
\]

Както при обединението, схемите трябва да бъдат съвместими. Резултатът съдържа всички кортежи, които са в $R$, но не са в $S$:

\[
R-S=\{t\mid t\in R \text{ и } t\notin S\}.
\]

Разликата не е комутативна. В общия случай

\[
R-S\neq S-R.
\]

Редът на операндите има значение: първата релация определя множеството, от което се изваждат кортежи.

\subsection{Сечение}

Сечението на $R$ и $S$ се означава с

\[
R\cap S.
\]

Операцията се прилага върху релации със съвместими схеми и връща релация, съдържаща само кортежите, общи за двете релации:

\[
R\cap S=\{t\mid t\in R \text{ и } t\in S\}.
\]

Сечението е комутативно и асоциативно:

\[
R\cap S=S\cap R,
\]

\[
(R\cap S)\cap T=R\cap(S\cap T).
\]

То може да бъде изразено чрез основните операции:

\[
R\cap S=R-(R-S).
\]

\subsection{Проекция}

Проекцията е унарна операция, която избира определени атрибути от релация. Означава се с

\[
\pi_{A_1,A_2,\ldots,A_k}(R),
\]

където $A_1,A_2,\ldots,A_k$ е списък от атрибути.

Резултатът съдържа само посочените колони. Ако след премахването на останалите колони се получат еднакви кортежи, повторенията се премахват. Например

\[
\pi_{\textit{Title},\textit{Year},\textit{Length}}(\textit{Movie})
\]

избира трите посочени атрибута от релацията \textit{Movie}.

Проекцията променя схемата, като намалява броя на атрибутите. Затова тя често се описва като вертикално отсяване на релацията.

\subsection{Селекция}

Селекцията, наричана още рестрикция, е унарна операция, която избира кортежите, удовлетворяващи определено условие. Означава се с

\[
\sigma_C(R),
\]

където $C$ е предикат.

Предикатът може да сравнява два атрибута или атрибут с константа. Допустимите оператори са

\[
=,\quad \neq,\quad <,\quad >,\quad \leq,\quad \geq.
\]

Пример:

\[
\sigma_{\textit{Year}<1980}(\textit{Movie}).
\]

Резултатът съдържа само филмите, чиято година е по-малка от $1980$. Схемата остава същата, но броят на кортежите може да намалее. Поради това селекцията се разглежда като хоризонтално отсяване на релацията.

Селекцията е комутативна при последователно прилагане:

\[
\sigma_{C_1}(\sigma_{C_2}(R))
=
\sigma_{C_2}(\sigma_{C_1}(R)).
\]

\subsection{Декартово произведение}

Декартовото произведение на $R$ и $S$ се означава с

\[
R\times S.
\]

Резултатът съдържа всички двойки от кортежи, при които първият кортеж е от $R$, а вторият е от $S$. Схемата на резултата е обединение на схемите на $R$ и $S$.

Ако $R$ има $m$ кортежа, а $S$ има $n$ кортежа, резултатът съдържа всички възможни комбинации между тях. При еднакви имена на атрибути се използва квалифицирано име от вида

\[
\textit{имеНаРелация}.\textit{имеНаАтрибут}.
\]

Например могат да се използват \textit{Movie.Year} и \textit{StarsIn.MovieYear}.

Декартовото произведение е бинарна операция, комутативна и асоциативна според разглеждането на съответните атрибути и подредбата на схемата.

\subsection{Тета-съединение}

Тета-съединението на $R$ и $S$ при условие $C$ се означава с

\[
R\bowtie_C S.
\]

То се получава чрез декартово произведение, последвано от селекция:

\[
R\bowtie_C S=\sigma_C(R\times S).
\]

Резултатът има схема, получена чрез обединяване на схемите на $R$ и $S$, и съдържа само онези двойки кортежи, които удовлетворяват условието $C$.

Ако условието включва само съвпадение между атрибути, тета-съединението се нарича еквисъединение. Например условието

\[
\textit{Movie.Year}=\textit{StarsIn.MovieYear}
\]

свързва кортежи, за които стойностите на съответните атрибути съвпадат.

\subsection{Естествено съединение}

Естественото съединение се означава с

\[
R\bowtie S.
\]

То свързва релациите по всички атрибути с еднакви имена. Резултатът съдържа кортежите от декартовото произведение, за които стойностите на едноименните атрибути съвпадат. Повтарящите се колони се премахват автоматично.

Естественото съединение може да бъде представено чрез тета-съединение и проекция:

\[
R\bowtie S
=
\pi_L\left(\sigma_C(R\times S)\right),
\]

където $C$ е условието за равенство по едноименните атрибути, а $L$ съдържа всички атрибути на $R$ и онези атрибути на $S$, които не присъстват в $R$.

Ако общите атрибути са $A_1,\ldots,A_n$, условието е от вида

\[
C=(R.A_1=S.A_1)\mathbin{\mathrm{AND}}\cdots
\mathbin{\mathrm{AND}}(R.A_n=S.A_n).
\]

\subsection{Преименуване}

Преименуването се означава с оператора $\rho$. То променя името на релацията и имената на нейните атрибути, без да променя кортежите.

Записът

\[
R_1:=\rho_{R_1(A_1,\ldots,A_n)}(R_2)
\]

означава, че $R_1$ е релация с атрибути $A_1,\ldots,A_n$ и със същите кортежи като $R_2$.

Преименуването е полезно, когато една и съща релация участва повече от веднъж в алгебричен израз или когато имената на атрибутите трябва да бъдат разграничени.

\subsection{Класификация и приоритет}

Основните операции на релационната алгебра са:

\begin{itemize}
    \item обединение;
    \item разлика;
    \item декартово произведение;
    \item проекция;
    \item селекция;
    \item преименуване.
\end{itemize}

Допълнителните операции са:

\begin{itemize}
    \item сечение;
    \item тета-съединение;
    \item естествено съединение.
\end{itemize}

Те могат да се получат от основните операции чрез следните зависимости:

\[
R\cap S=R-(R-S),
\]

\[
R\bowtie_C S=\sigma_C(R\times S),
\]

\[
R\bowtie S=\pi_L\left(\sigma_C(R\times S)\right).
\]

Приоритетът на операторите в алгебричните изрази е:

\begin{enumerate}
    \item унарни оператори -- селекция, проекция и преименуване;
    \item декартово произведение и съединение;
    \item сечение;
    \item обединение и разлика;
    \item изрази в скоби.
\end{enumerate}

При нееднозначност се използват скоби.

\section{Примерни задачи}

\subsection{Съставяне на SQL заявки}

Нека е дадена таблицата \texttt{movies} със следните колони: \texttt{id}, \texttt{title}, \texttt{length} и \texttt{release\_date}.

Задача: да се изведат заглавията на филмите с ненулева дължина, по-малка от $120$, подредени по дата на издаване в низходящ ред.

\begin{verbatim}
SELECT title FROM movies
WHERE length IS NOT NULL AND length < 120
ORDER BY release_date DESC
\end{verbatim}

Задача: да се добави филм.

\begin{verbatim}
INSERT INTO movies (title, length, release_date)
VALUES ('Star wars', 124, '09-04-1984')
\end{verbatim}

Задача: да се промени заглавието на филма.

\begin{verbatim}
UPDATE movies
SET title = 'Star Wars'
WHERE title = 'Star wars'
\end{verbatim}

Задача: да се изтрият филмите с дължина под $120$.

\begin{verbatim}
DELETE FROM movies
WHERE length < 120
\end{verbatim}

След промени чрез \texttt{INSERT}, \texttt{UPDATE} или \texttt{DELETE} се изпълнява:

\begin{verbatim}
COMMIT
\end{verbatim}

\subsection{Пример за DDL}

\begin{verbatim}
CREATE TABLE movies (
  id INTEGER PRIMARY KEY,
  title CHAR(50) NOT NULL,
  length INTEGER NULL
);
\end{verbatim}

Добавяне на атрибут:

\begin{verbatim}
ALTER TABLE movies ADD release_date DATE;
\end{verbatim}

Премахване на атрибут:

\begin{verbatim}
ALTER TABLE movies DROP COLUMN length;
\end{verbatim}

Изтриване на таблицата:

\begin{verbatim}
DROP TABLE movies;
\end{verbatim}

\subsection{Тригери}

Тригерите са свързани с автоматично изпълнение на действие при настъпване на определено събитие върху данни. В предоставения материал не са зададени конкретен синтаксис, събития или примерна реализация на тригер. Поради това тук не се фиксира конкретен SQL вариант за тригер, тъй като синтаксисът зависи от използваната система за управление на бази от данни.

\section{Изпитен фокус}

За устно изложение трябва да могат да бъдат възпроизведени:

\begin{itemize}
    \item определенията за домейн, релация, атрибут, кортеж, схема и екземпляр;
    \item разликата между схема на релация и текущ екземпляр на релация;
    \item определението за схема на релационна база от данни;
    \item изискването компонентите на кортежите да бъдат атомарни;
    \item стъпките при проектиране и реализиране на релационна база от данни;
    \item предназначението на DDL и командите \texttt{CREATE}, \texttt{ALTER}, \texttt{DROP};
    \item предназначението на DML и командите \texttt{SELECT}, \texttt{INSERT}, \texttt{UPDATE}, \texttt{DELETE};
    \item ролята на \texttt{COMMIT} след DML операции;
    \item дефинициите и условията за приложимост на обединение, разлика и сечение;
    \item разликата между селекция и проекция;
    \item определението за декартово произведение;
    \item изразяването на тета-съединение чрез декартово произведение и селекция;
    \item особеността на естественото съединение да свързва по едноименни атрибути и да премахва повтарящите се колони;
    \item ролята на преименуването;
    \item класификацията на основни и допълнителни операции;
    \item зависимостите
    \[
    R\cap S=R-(R-S),\qquad
    R\bowtie_C S=\sigma_C(R\times S);
    \]
    \item приоритетът на операторите в релационната алгебра;
    \item съставянето на примерни SQL заявки за извличане, добавяне, промяна и изтриване на данни;
    \item примерите за DDL операции и предназначението на тригерите.
\end{itemize}
